\chapter{Summary}
\paragraph{Scope of the summarized results.} The results summarized below originate from two complementary studies. The canonical study (Chapters~\ref{sec:experiments}-- \ref{sec:extensions}) evaluates relay processing on a memoryless channel with white noise and uncoded transmission. The unknown-channel study (Chapter~\ref{sec:unknown-channels}) extends that baseline by deliberately introducing impairments that violate those assumptions, including intersymbol interference (ISI), nonlinear distortion, and model mismatch. Conclusions involving sequence detection, Viterbi MLSE, pilot overhead, or blind equalization therefore refer specifically to that extension rather than to the canonical system model.

This work examined classical and neural network–based relay strategies for \textit{two-hop relay communication} using a unified simulation framework. The objective was to analyze performance–complexity tradeoffs and determine when learning-based relays provide practical advantages over analytically derived methods.
The entire study is carried out on one canonical setup, fixed in advance and never departed from: a two-hop SISO link over i.i.d.\ Rayleigh fast fading, complex baseband, QPSK, uncoded BER, with the relay function as the only varied axis. The study showed the following:
\REV{``two-hop cooperative communication'' \(\to\) ``two-hop relay communication'', consistent with the Introduction and Abstract: the studied system has no direct source--destination link, so ``cooperative'' (which conventionally implies a direct path combined with the relayed path) is replaced by ``relay''.}
\begin{enumerate}
    \item Neural relays yield a consistent performance improvement over amplify-and-forward (AF) at low signal-to-noise ratio (SNR), where their learned soft denoising avoids AF's noise amplification (H1).
    \item Classical \textit{decode-and-forward (DF)} matched or outperformed all neural approaches at \textbf{medium-to-high SNR} while requiring no trainable parameters, confirming that DF attains the lowest per-symbol BER in this regime (H2)..
    \item A systematic complexity analysis demonstrated that performance \textbf{does not improve monotonically with model size} (H3).
    \item A \textbf{minimal neural relay} (169 parameters) achieved performance comparable to significantly larger models, while intermediate-sized networks exhibited overfitting.
    \item When neural relays were evaluated under equal parameter budgets, architectural differences largely vanished, indicating that \textbf{model capacity}, rather than architectural class, dominates performance for this task (H4).

    \item Sequence-based neural relays such as structured state-space models achieved comparable error performance with improved training efficiency (H5).
    \item From a deployment perspective, a \texttt{hybrid relay} combining neural processing at low SNR with DF at high SNemerged as the \textbf{recommended default}, achieving near-optimal performance across the SNR range with minimal overhead
\end{enumerate}    
    The first seven findings concern the canonical memoryless relay model; the remaining findings summarize the deliberate model-mismatch extension of Chapter~\ref{sec:unknown-channels}.
\begin{enumerate}
\setcounter{enumi}{7}
    \item Beyond the canonical memoryless model, on an \textbf{unknown channel} violating classical assumptions (unmodeled intersymbol interference or biased nonlinearity), AF and DF failed to relay reliably (exhibiting an analytic BER floor of $0.25$, with DF \emph{worsening} as SNR increased) while the same minimal 169-parameter MLP restored reliable relaying by learning the impairment from data, including an arbitrary composite cascade of ISI, amplifier nonlinearity, and unknown phase. It tracked $\approx$1--1.5 dB behind a \emph{pilot-aided} Viterbi receiver on pure ISI ($\approx$2 dB on the composite cascade) and, in the fully posterior-free regime (no pilots, no channel prior), matched the best blind classical equalizer while avoiding its instability and online-adaptation cost (H6).

    \item Sweeping the intermediate \emph{partial}-posterior case located the crossover precisely: it is governed by the \textbf{reliability of the channel estimate} and by pilot overhead. Pilot-aided classical processing wins on long blocks at negligible cost, but the zero-pilot learned relay becomes decisive once blocks are too short, or channels too transient, for pilots to yield an accurate estimate. The learned relay's advantage is therefore specific to \emph{structural} model-class mismatch (memory, nonlinearity, absent pilots) rather than to unknownness per se. Together with the high-SNR result, this delineates precisely where learned relays add value: never above the model-aware optimum, but uniquely robust to \emph{not knowing which model applies} and to the pilot overhead of channel identification.

    \item A \textbf{complexity accounting} completed the trade-off: the learned relay's per-symbol cost is constant in modulation order and channel memory, and $30$--$90\times$ faster in measured wall-clock as a vectorized operation for the reference implementations compared (Section~\ref{sec:complexity-viterbi-mlp}), whereas exact Viterbi MLSE grows as $M^{L}$ and becomes intractable. The learned relay thus trades a small BER gap for cost that scales gracefully exactly where classical sequence detection does not.

    \item In the modulation extension, the canonical findings generalized fully to QPSK via I/Q independence, while 16-QAM exposed a structural error floor of per-axis processing; reformulating the relay as a \textbf{joint $N$-class 2D classifier} eliminated the floor and matched DF at high SNR.
\end{enumerate}

Overall, the results indicate that learning-based relays are best used as targeted enhancements rather than universal replacements, with simplicity, parameter efficiency, and hybrid design as key principles for practical relay systems. All conclusions hold for the \emph{uncoded}, symbol-level canonical setting studied here (Remark~\ref{rem:df-terminology}); extending the comparison to coded block-DF, soft-information (LLR) forwarding relays, other channel families, multi-antenna configurations, and denser constellations (64-QAM and beyond) is identified in Chapter~\ref{sec:discussion-and-conclusions} as the principal direction for future work.
\REV{Fix applied (structure and precision): the unknown-channel summary was a single $\sim$250-word item that buried the pilot-crossover and complexity results; it is now split into the three items above, mirroring the structure of Chapter~\ref{sec:unknown-channels}. ``A few dB behind a pilot-aided Viterbi receiver'' was replaced by the precise figures used in the body ($\approx$1--1.5~dB on pure ISI, $\approx$2~dB on the composite), and the repeated $30$--$90\times$ wall-clock figure now carries the same implementation caveat as Section~\ref{sec:complexity-viterbi-mlp}.}
